%% 05-the-analytics-engine.tex -- generated by tools/build_from_docx.py; edit freely

\hypertarget{the-analytics-engine}{%
\section{The Analytics Engine}\label{the-analytics-engine}}

The Analytics Engine is the capability through which an agent draws
conclusions from the evidence its Monitoring System produced. It
extracts features from observation streams and from the Data Continuum,
performs inference, and publishes insight: anomaly indications,
forecasts, classifications, derived key performance indicators and
root-cause hypotheses, each carrying the uncertainty of the model that
produced it and the quality of the evidence it consumed.

The Analytics Engine is where contributed algorithms predominantly live.
Forecasting models, energy-optimization models, semantic analysis,
failure prediction and the analytical half of routing algorithms are all
Analytics Engine work, and the framework constrains none of them in how
they reason. What it constrains is what they may consume, what they may
emit, and what the platform establishes about their output before
anything downstream relies on it.

One boundary governs the chapter, and it is what makes a contributed
model substitutable. The Analytics Engine produces insight and never
action. It holds no authority over any managed object, proposes no
change to one, and names no controller. A partner may therefore replace
a model, change its architecture or retrain it on different data without
any change to what the agent is permitted to do, because the Analytics
Engine was never what permitted it.

\hypertarget{role-and-scope-of-the-analytics-engine}{%
\subsection{Role and Scope of the Analytics Engine}\label{role-and-scope-of-the-analytics-engine}}

\hypertarget{position-in-the-loop-1}{%
\subsubsection{Position in the loop}\label{position-in-the-loop-1}}

The Analytics Engine occupies the second position of the loop. It
consumes observations from the Monitoring System of its own agent over
the local runtime binding, or insight from a provider where it is
delegated, or data products from the Data Continuum where its features
are drawn from derived artifacts. It publishes insight to the Decision
Engine of its own agent, and, where its descriptor exposes a contract,
to other agents as a data product or through an exposed contract.

The canonical deployment splits the Monitoring System at the edge from
the Analytics Engine at a regional or central tier, which is also where
declared graphics processing resources are consumed. That split is
permitted only where the resulting effective latency class, computed at
admission from the declared distribution and the declared class of each
transport hop, still meets what the agent's loop requires.

\hypertarget{capability-declarations-1}{%
\subsubsection{Capability declarations}\label{capability-declarations-1}}

At \texttt{full} the Analytics Engine performs its analytics end to end.
At \texttt{partial} it shall enumerate the insight types it supports; an
unenumerated partial declaration is rejected at admission, and the
enumeration is thereafter enforced at the publication boundary rather
than trusted. At \texttt{delegated} it consumes insight produced
elsewhere, naming the provider and the interface. At \texttt{none} the
agent acts on externally supplied insight or on policy alone.

The declaration determines what the agent may publish, not how well it
reasons. An Analytics Engine that publishes an insight type it did not
declare has its insight dropped at the boundary and a conformance fault
recorded against the agent.

\hypertarget{input-quality-and-model-confidence}{%
\subsection{Input Quality and Model Confidence}\label{input-quality-and-model-confidence}}

\hypertarget{two-distinct-quantities}{%
\subsubsection{Two distinct quantities}\label{two-distinct-quantities}}

Every insight carries two quality structures, and they are never merged.
The model confidence is the model's own expression of uncertainty about
its output, produced by the contributed logic. The input quality is a
statement about the evidence the inference consumed, computed by the
Analytics Engine runtime from the envelopes of the observations actually
used: the minimum coverage fraction among them, the maximum staleness
among them, the identifiers of the contributing observation envelopes,
and a derived flag recording whether the input was degraded.

The reason for the separation is that a model cannot report what it
cannot know. A model fed telemetry with substantial gaps reports
ordinary confidence, because its confidence expresses uncertainty about
the relationship it was trained on, not about whether the evidence it
received was trustworthy. Nothing inside the model can distinguish a
quiet network from a network whose observations did not arrive. The
platform can, because the Monitoring System computed coverage against a
resolved scope, and the Analytics Engine runtime reads that computation
directly from the envelopes consumed.

Input quality is therefore never authored by contributed logic, and an
insight published without it is rejected at the publication boundary.

\hypertarget{downstream-consequence}{%
\subsubsection{Downstream consequence}\label{downstream-consequence}}

The Decision Engine evaluates input quality independently of model
confidence. Where the degraded-input flag is set, the Decision Engine's
declared degraded-evidence behavior applies regardless of how confident
the model claims to be. A high-confidence insight drawn from degraded
evidence is thereby prevented from producing a decision that its
evidence never supported.

This is what makes degradation propagate along the loop by construction
rather than by a model noticing its own predicament. Coverage computed
at the Monitoring System reaches the Decision Engine as a
decision-relevant property, having passed through the Analytics Engine
without depending on the model's cooperation at any point.

\hypertarget{validity-containment-and-authority}{%
\subsection{Validity, Containment and Authority}\label{validity-containment-and-authority}}

\hypertarget{validity-conditions-are-machine-evaluable}{%
\subsubsection{Validity conditions are machine-evaluable}\label{validity-conditions-are-machine-evaluable}}

Every model declares the conditions under which its output is valid, and
those conditions shall be machine-evaluable predicates over observable
input features. The predicates follow the structured form defined for
the framework: comparisons and combinators over named features, with no
expression language.

The Analytics Engine runtime evaluates them against the consumed
observations before inference, and stamps each insight with a validity
state of in-envelope or out-of-envelope. The stamp is applied by the
runtime and not by the model, so that a model operating outside the
conditions its own publisher declared is visible as such to everything
downstream, whatever the model reports about itself.

\hypertarget{containment-rather-than-verification}{%
\subsubsection{Containment rather than verification}\label{containment-rather-than-verification}}

The platform does not attempt to verify a publisher's validity claim. It
has no ground truth against which to do so and cannot acquire one, and a
verification the platform cannot perform is not made sound by asserting
it. The framework substitutes containment: the claim is not checked, but
its consequences are bounded, observed and acted on.

Containment has four parts. The runtime stamps validity state on every
insight, so that an out-of-envelope inference is never mistaken for an
ordinary one. An evaluation function tracks outcome against declared
objective for each agent independently of what the agent reports.
Sustained divergence between outcome and objective triggers an automatic
downgrade of the agent's authority, against a declared and auditable
threshold rather than a discretionary judgment. And authority once
withdrawn is restored only by an explicit review that records its
outcome.

The asymmetry between the two directions is deliberate. Downgrade is
automatic because containment that waits is not containment: an agent
diverging from its objective is narrowed and then suspended without
anyone being consulted. Restoration is not automatic, and in particular
is not the expiry of a timer: the downgrade does not lapse, and no
interval of good behavior restores authority by itself. A review
decides, and the decision is recorded. The reviewer need not be a person
--- the evaluation function may restore authority on declared evidence,
such as sustained conformance across a declared window --- but
restoration is in every case a decided act with an auditable record
behind it, never a silence into which authority returns.

\hypertarget{the-separation-of-two-authorities}{%
\subsubsection{The separation of two authorities}\label{the-separation-of-two-authorities}}

The evaluation function does not veto individual decisions. A
per-decision veto would constitute a second reasoning entity in the
loop, which the framework does not admit. The evaluation function holds
lifecycle authority instead: it may narrow an agent's scope, suspend its
emission of decisions, or withdraw it, and it does so between cycles
rather than within one.

Per-decision authority remains with the conflict fabric. The two
authorities never contend, because they act on different objects at
different time scales: one governs a single decision as it is taken, the
other governs an agent across many. A downgrade is enacted by the domain
management function as a registry change and applied by the platform at
the intent egress gate, and the agent cannot decline it.

\hypertarget{model-governance}{%
\subsection{Model Governance}\label{model-governance}}

Each model an Analytics Engine uses is declared with its identifier and
version, the feature set it requires and the version of that feature
set, its validity conditions, and its serving mode. A model is bound by
pinned identity and digest and never by a mutable reference such as a
tag denoting the latest artifact.

The serving mode is declared as either an artifact embedded in the
capability or an external endpoint, and the two are contractually
equivalent: both produce the same insight schema, and an endpoint
additionally declares its measured inference latency. The declaration of
serving mode follows from the latency class, since a streaming inference
endpoint and a batch job are not interchangeable within a loop that has
committed to a deadline.

Model lifecycle operations --- training, retraining, recalibration,
validation, promotion and interaction with the model registry --- never
occur on the decision path. A model that is being retrained is not a
model that is serving, and the framework admits no construction in which
a cycle waits on a training operation.

Inference latency is a mandatory observability metric emitted by the
runtime rather than a figure declared by the publisher. Sustained breach
of the agent's effective latency class is a conformance fault, feeds the
downgrade path of Section 5.3, and causes the effective class to be
re-stamped downward in the registry so that composition and priority
logic use the measured figure rather than the intended one.

\hypertarget{the-insight-contract-and-envelope}{%
\subsection{The Insight Contract and Envelope}\label{the-insight-contract-and-envelope}}

\hypertarget{insight-schemas-carry-no-actions}{%
\subsubsection{Insight schemas carry no actions}\label{insight-schemas-carry-no-actions}}

An insight schema shall not contain action identifiers, target
configurations, parameter values to be applied, or controller
references. Admission rejects any insight contract containing
actuation-shaped fields.

The prohibition is structural rather than stylistic. An action expressed
as an insight would carry no action-space bounds, no conflict-resolution
policy and no authority envelope, because the Analytics Engine has none
to give it. It would be decision authority laundered through a
capability that holds none, and everything the framework establishes
about what an agent may change would be bypassed by the simple expedient
of naming a change as an observation about the world.

Root-cause hypotheses remain permitted and are not affected by the
prohibition. A hypothesis names a cause; it does not name a remedy. The
distinction is precisely the boundary between the Analytics Engine and
the Decision Engine, and it is enforced at the schema.

\hypertarget{provenance}{%
\subsubsection{Provenance}\label{provenance}}

Every insight envelope carries the identifier and version of the
producing agent together with the identifier and version of the model
that produced it. A consumer of delegated analytics declares the
provider versions it accepts as a pinned range.

A provider model or agent version outside the consumer's accepted range
raises a binding fault at the consumer and triggers the consumer's
declared behavior for an unavailable model. It does not silently alter
the consumer's decision behavior. Without this, a provider-side model
upgrade would change the behavior of every consuming agent with no
change to any consumer's own descriptor or version --- drift across an
ownership boundary that nothing downstream could detect.

\hypertarget{declared-failure-behavior-1}{%
\subsection{Declared Failure Behavior}\label{declared-failure-behavior-1}}

The Analytics Engine's failure policy declares the five responses of
Table 5-1.

It declares the response to stale input, as emission of an
unavailable-insight signal, degradation of reported confidence, or a
declared fallback. It declares the minimum input coverage fraction below
which the capability shall emit an unavailable-insight signal rather
than infer. It declares the response to operation outside model
validity, as either an unavailable-insight signal or emission carrying
the validity flag, silent extrapolation being prohibited in every case.
It declares the behavior on loss of a serving endpoint. And it declares
the response to drift, which shall at minimum emit a drift signal.

A hard invariant governs fallbacks: a fallback heuristic may be used
only where it was declared in advance. An undeclared fallback is
indistinguishable downstream from the model whose place it takes, and
would change the agent's behavior without changing its contract ---
which is exactly the condition the framework exists to prevent.

\begin{wptable}{5-1}

\begin{longtable}[]{@{}
  >{\raggedright\arraybackslash}p{(\columnwidth - 6\tabcolsep) * \real{0.2625}}
  >{\raggedright\arraybackslash}p{(\columnwidth - 6\tabcolsep) * \real{0.1250}}
  >{\raggedright\arraybackslash}p{(\columnwidth - 6\tabcolsep) * \real{0.1000}}
  >{\raggedright\arraybackslash}p{(\columnwidth - 6\tabcolsep) * \real{0.5125}}@{}}
\caption{Analytics Engine failure policy keys}\tabularnewline
\toprule\noalign{}
\begin{minipage}[b]{\linewidth}\raggedright
\textbf{Field in failure\_\allowbreak{}policy.\allowbreak{}AE}
\end{minipage} & \begin{minipage}[b]{\linewidth}\raggedright
\textbf{Type}
\end{minipage} & \begin{minipage}[b]{\linewidth}\raggedright
\textbf{Req.}
\end{minipage} & \begin{minipage}[b]{\linewidth}\raggedright
\textbf{Definition}
\end{minipage} \\
\midrule\noalign{}
\endfirsthead
\toprule\noalign{}
\begin{minipage}[b]{\linewidth}\raggedright
\textbf{Field in failure\_\allowbreak{}policy.\allowbreak{}AE}
\end{minipage} & \begin{minipage}[b]{\linewidth}\raggedright
\textbf{Type}
\end{minipage} & \begin{minipage}[b]{\linewidth}\raggedright
\textbf{Req.}
\end{minipage} & \begin{minipage}[b]{\linewidth}\raggedright
\textbf{Definition}
\end{minipage} \\
\midrule\noalign{}
\endhead
\bottomrule\noalign{}
\endlastfoot
\textbf{input\_\allowbreak{}staleness} & enum & M &
insight\_\allowbreak{}unavailable, degrade\_\allowbreak{}confidence or
declared\_\allowbreak{}fallback.\allowbreak{} \\
\textbf{min\_\allowbreak{}input\_\allowbreak{}coverage} & number & M &
Coverage fraction below which insight\_\allowbreak{}unavailable shall be
emitted. \\
\textbf{model\_\allowbreak{}out\_\allowbreak{}of\_\allowbreak{}validity}
& enum & M & insight\_\allowbreak{}unavailable or
emit\_\allowbreak{}with\_\allowbreak{}validity\_\allowbreak{}flag.\allowbreak{}
Silent extrapolation is prohibited. \\
\textbf{model\_\allowbreak{}unavailable} & enum & M & Behavior on loss
of a serving endpoint. \\
\textbf{drift\_\allowbreak{}response} & enum & M & Shall at minimum emit
a drift signal. \\
\end{longtable}

\end{wptable}

\hypertarget{the-insight-contract}{%
\subsection{The Insight Contract}\label{the-insight-contract}}

An insight contract declares what an Analytics Engine publishes, over
which subject classes, from which inputs, under which models and
validity conditions, and at what freshness. Its fields are given in
Table 5-2. A consumer binds against the contract identifier and a pinned
compatibility range, so that a change in what is published reaches the
consumer as a binding fault rather than as a silent change in meaning.

\begin{wptable}{5-2}

\begin{longtable}[]{@{}
  >{\raggedright\arraybackslash}p{(\columnwidth - 6\tabcolsep) * \real{0.2625}}
  >{\raggedright\arraybackslash}p{(\columnwidth - 6\tabcolsep) * \real{0.1250}}
  >{\raggedright\arraybackslash}p{(\columnwidth - 6\tabcolsep) * \real{0.1499}}
  >{\raggedright\arraybackslash}p{(\columnwidth - 6\tabcolsep) * \real{0.4625}}@{}}
\caption{Insight contract fields}\tabularnewline
\toprule\noalign{}
\begin{minipage}[b]{\linewidth}\raggedright
\textbf{Field in insight\_\allowbreak{}contract}
\end{minipage} & \begin{minipage}[b]{\linewidth}\raggedright
\textbf{Type}
\end{minipage} & \begin{minipage}[b]{\linewidth}\raggedright
\textbf{Req.}
\end{minipage} & \begin{minipage}[b]{\linewidth}\raggedright
\textbf{Definition}
\end{minipage} \\
\midrule\noalign{}
\endfirsthead
\toprule\noalign{}
\begin{minipage}[b]{\linewidth}\raggedright
\textbf{Field in insight\_\allowbreak{}contract}
\end{minipage} & \begin{minipage}[b]{\linewidth}\raggedright
\textbf{Type}
\end{minipage} & \begin{minipage}[b]{\linewidth}\raggedright
\textbf{Req.}
\end{minipage} & \begin{minipage}[b]{\linewidth}\raggedright
\textbf{Definition}
\end{minipage} \\
\midrule\noalign{}
\endhead
\bottomrule\noalign{}
\endlastfoot
\textbf{contract\_\allowbreak{}id} & string & M & Identifier of the
insight contract. \\
\textbf{version} & string & M & Version of the contract; consumers pin a
compatibility range. \\
\textbf{insight\_\allowbreak{}types{[}{]}} & array & M & Types
published. A partial AE shall enumerate these exhaustively. \\
\textbf{subject\_\allowbreak{}classes{[}{]}} & array & M & Classes of
managed object the insights concern. \\
\textbf{schema\_\allowbreak{}ref} & string & M & Semantic schema of the
payload, with its version. Rejected at admission if it carries
actuation-shaped fields. \\
\textbf{inputs{[}{]}} & array & M & Observation contracts or data
products consumed. \\
\textbf{models{[}{]}} & array & C: model-produced & Declared models with
identifier, version, feature set and serving mode. \\
\textbf{validity\_\allowbreak{}conditions} & object & C: model-produced
& Machine-evaluable predicates over observable input features. \\
\textbf{freshness} & duration & M & Maximum age at which a consumer may
rely on the insight. \\
\textbf{binding\_\allowbreak{}classes{[}{]}} & array & M & Permitted
bindings: LOCAL\_\allowbreak{}RUNTIME, DATA\_\allowbreak{}CONTINUUM or
SMB. \\
\end{longtable}

\end{wptable}

Table A-1 gives the descriptor location at which each field of Tables
5-1 and 5-2 is declared.

Unlike an observation contract, an insight contract may be carried on
the service-based management bus as well as over the local runtime
binding and the Data Continuum, because insight is control-plane content
rather than telemetry. A contract whose schema carries an action
identifier, a target configuration, a parameter value to be applied or a
controller reference is rejected at admission, whatever binding it
declares.

\hypertarget{the-insight-envelope}{%
\subsection{The Insight Envelope}\label{the-insight-envelope}}

Every insight carries the envelope of Table 5-3, populated by the
Analytics Engine runtime. An insight reaching a publication boundary
without a complete envelope is rejected there.

\begin{wptable}{5-3}

\begin{longtable}[]{@{}
  >{\raggedright\arraybackslash}p{(\columnwidth - 6\tabcolsep) * \real{0.2625}}
  >{\raggedright\arraybackslash}p{(\columnwidth - 6\tabcolsep) * \real{0.1250}}
  >{\raggedright\arraybackslash}p{(\columnwidth - 6\tabcolsep) * \real{0.1499}}
  >{\raggedright\arraybackslash}p{(\columnwidth - 6\tabcolsep) * \real{0.4625}}@{}}
\caption{Insight envelope fields}\tabularnewline
\toprule\noalign{}
\begin{minipage}[b]{\linewidth}\raggedright
\textbf{Field in envelope}
\end{minipage} & \begin{minipage}[b]{\linewidth}\raggedright
\textbf{Type}
\end{minipage} & \begin{minipage}[b]{\linewidth}\raggedright
\textbf{Req.}
\end{minipage} & \begin{minipage}[b]{\linewidth}\raggedright
\textbf{Definition}
\end{minipage} \\
\midrule\noalign{}
\endfirsthead
\toprule\noalign{}
\begin{minipage}[b]{\linewidth}\raggedright
\textbf{Field in envelope}
\end{minipage} & \begin{minipage}[b]{\linewidth}\raggedright
\textbf{Type}
\end{minipage} & \begin{minipage}[b]{\linewidth}\raggedright
\textbf{Req.}
\end{minipage} & \begin{minipage}[b]{\linewidth}\raggedright
\textbf{Definition}
\end{minipage} \\
\midrule\noalign{}
\endhead
\bottomrule\noalign{}
\endlastfoot
\textbf{insight\_\allowbreak{}id} & uuid & M & Unique identifier of the
insight, assigned by the AE runtime. \\
\textbf{agent\_\allowbreak{}instance\_\allowbreak{}id} & string & M &
Instance that produced the insight. \\
\textbf{agent\_\allowbreak{}version} & string & M & Version of the
producing agent. \\
\textbf{contract\_\allowbreak{}id} & string & M & Insight contract and
version under which the insight is published. \\
\textbf{insight\_\allowbreak{}type} & enum & M & anomaly, forecast,
classification, derived\_\allowbreak{}kpi or
root\_\allowbreak{}cause\_\allowbreak{}hypothesis.\allowbreak{} \\
\textbf{subject} & string & M & Managed object or scope selector the
insight concerns. \\
\textbf{payload} & object & M & Content under its registered semantic
schema, carrying no actuation-shaped field. \\
\textbf{model\_\allowbreak{}id} & string & C: model-produced &
Identifier of the model that produced the insight. \\
\textbf{model\_\allowbreak{}version} & string & C: model-produced &
Version of that model. \\
\textbf{model\_\allowbreak{}confidence} & number & C: declared by type &
The model\textquotesingle s own uncertainty output. Never a substitute
for input quality. \\
\textbf{input\_\allowbreak{}quality} & object & M & Minimum input
coverage, maximum input staleness, contributing observation envelope
identifiers and the derived degraded\_\allowbreak{}input flag. Computed
by the runtime. \\
\textbf{validity\_\allowbreak{}state} & enum & M &
in\_\allowbreak{}envelope or
out\_\allowbreak{}of\_\allowbreak{}envelope, stamped by the runtime from
the declared validity predicates. \\
\textbf{fallback\_\allowbreak{}used} & bool & M & Whether a declared
fallback produced the output in place of the model. \\
\textbf{inference\_\allowbreak{}latency} & duration & M & Measured
inference latency, emitted by the runtime. \\
\textbf{event\_\allowbreak{}time} & timestamp & M & Time to which the
insight refers. \\
\textbf{produced\_\allowbreak{}time} & timestamp & M & Time the insight
was produced. \\
\textbf{feature\_\allowbreak{}set\_\allowbreak{}version} & string & C:
model-produced & Version of the feature set consumed. \\
\end{longtable}

\end{wptable}

Four envelope fields carry the chapter's substance. Input quality and
model confidence stand side by side and are never reconciled into a
single number, because they answer different questions. The validity
state records whether the runtime found the input inside the envelope
the publisher declared. The fallback flag makes a declared fallback
visible as such, so that a consumer knows the output did not come from
the model it expected. And the measured inference latency is what the
registry uses to re-stamp an effective latency class, rather than the
figure the descriptor declared.

\hypertarget{delegated-analytics}{%
\subsection{Delegated Analytics}\label{delegated-analytics}}

An analytics delegation may cross domain boundaries, as a data product
or as an exposed contract, which distinguishes it from monitoring and
actuation delegations. The reason is that insight carries no authority:
consuming a forecast produced in another domain grants the consumer
nothing over that domain, whereas consuming an actuation capability
would. An agent of the radio access domain that consumes a transport
domain's degradation forecast learns something it could not have
observed; its action space is unchanged, and no transport parameter
becomes available to it.

A distinction of level follows. An agent that reasons across several
domains at once is a global agent operating at the cross-domain
management level, which holds that breadth by construction; a domain
agent consuming another domain's insight remains a domain agent, bounded
by its own action space. Decisions may be exchanged between domain
agents, and such an exchange is a proposal: a decision received from a
peer binds nothing and confers no authority over the receiving agent,
which remains answerable for what it applies within its own scope. The
final decision is taken by the overlooking agent at the cross-domain
management level. Authority therefore descends, while knowledge moves
freely in every direction.

Three conditions apply. The delegation pins the provider's interface
profile together with a range of accepted provider and model versions.
Admission verifies that the provider's declared subject classes cover
the consumer's requirements and that the provider's declared freshness
meets the consumer's assumptions. And the provider's input quality and
validity state are forwarded to the consumer within the insight itself,
so that a consumer of delegated analytics evaluates the evidence behind
an insight on the same terms as an agent that produced it locally.

\hypertarget{requirements-2}{%
\subsection{Requirements}\label{requirements-2}}

The requirements of this section apply to every agent that declares an
Analytics Engine in any mode other than none, and to the platform that
admits and runs it.

\begin{requirement}{Req-T23-AE-01}{MUST}{Insight, never action}
\reqfield{Description / rationale}{%
The AE shall produce insight only. An insight schema shall not contain
action identifiers, target configurations, parameter values to be
applied or controller references, and admission shall reject any insight
contract carrying actuation-shaped fields. Root-cause hypotheses remain
permitted: a hypothesis names a cause, not a remedy.}
\reqfield{Dependency}{%
Req-T23-FW-02}
\reqfield{Traceability (Forward)}{%
\begin{itemize}
\tightlist
\item
  Checked at admission against the registered insight schema.
\item
  Enforced at the publication boundary.
\end{itemize}}
\reqfield{UNITY-6G components}{%
\begin{itemize}
\tightlist
\item
  Tuple descriptor
\item
  AE runtime
\item
  Admission (IDMO/DMO)
\end{itemize}}
\reqfield{Linked KPIs}{%
\begin{itemize}
\tightlist
\item
  KPI\#1: Deliver four distributed management entities (MS, AE, DE and
  ACT) deployable in each heterogeneous domain.
\item
  KPI\#42: DT-validated AI algorithms that do not break the
  time-sensitive operation when deployed in 99\% of the outputs.
\end{itemize}}
\end{requirement}

\begin{requirement}{Req-T23-AE-02}{MUST}{Input quality computed by the runtime}
\reqfield{Description / rationale}{%
Every insight shall carry an input quality structure reporting minimum
input coverage, maximum input staleness, the contributing observation
envelope identifiers and a derived degraded-input flag. It shall be
computed by the AE runtime from the envelopes actually consumed and
shall never be authored by contributed logic. An insight published
without it shall be rejected at the publication boundary.}
\reqfield{Dependency}{%
Req-T23-MS-07}
\reqfield{Traceability (Forward)}{%
\begin{itemize}
\tightlist
\item
  Computed by the AE runtime from consumed observation envelopes.
\end{itemize}}
\reqfield{UNITY-6G components}{%
\begin{itemize}
\tightlist
\item
  AE runtime
\item
  MS runtime
\end{itemize}}
\reqfield{Linked KPIs}{%
\begin{itemize}
\tightlist
\item
  KPI\#42: DT-validated AI algorithms that do not break the
  time-sensitive operation when deployed in 99\% of the outputs.
\item
  KPI\#36: Guarantee 99.99999\% service availability and reliability
  based on effective monitoring and LCM.
\end{itemize}}
\end{requirement}

\begin{requirement}{Req-T23-AE-03}{MUST}{Input quality evaluated independently of confidence}
\reqfield{Description / rationale}{%
The DE shall evaluate input quality independently of model confidence,
and a set degraded-input flag shall trigger the DE\textquotesingle s
declared degraded-evidence behavior regardless of the confidence the
model reports.}
\reqfield{Dependency}{%
Req-T23-AE-02}
\reqfield{Traceability (Forward)}{%
\begin{itemize}
\tightlist
\item
  Enforced by the DE runtime at policy evaluation.
\end{itemize}}
\reqfield{UNITY-6G components}{%
\begin{itemize}
\tightlist
\item
  DE runtime
\item
  AE runtime
\end{itemize}}
\reqfield{Linked KPIs}{%
\begin{itemize}
\tightlist
\item
  KPI\#42: DT-validated AI algorithms that do not break the
  time-sensitive operation when deployed in 99\% of the outputs.
\item
  KPI\#36: Guarantee 99.99999\% service availability and reliability
  based on effective monitoring and LCM.
\end{itemize}}
\end{requirement}

\begin{requirement}{Req-T23-AE-04}{MUST}{Machine-evaluable validity conditions}
\reqfield{Description / rationale}{%
Declared model validity conditions shall be machine-evaluable predicates
over observable input features, expressed in the structured predicate
form of the framework. The AE runtime shall evaluate them against the
consumed observations before inference and stamp each insight with a
validity state of in\_\allowbreak{}envelope or
out\_\allowbreak{}of\_\allowbreak{}envelope.\allowbreak{} The stamp
shall not be applied by contributed logic.}
\reqfield{Dependency}{%
Req-T23-AE-02}
\reqfield{Traceability (Forward)}{%
\begin{itemize}
\tightlist
\item
  Evaluated and stamped by the AE runtime.
\item
  Checked for evaluability at admission.
\end{itemize}}
\reqfield{UNITY-6G components}{%
\begin{itemize}
\tightlist
\item
  AE runtime
\item
  Tuple descriptor
\end{itemize}}
\reqfield{Linked KPIs}{%
\begin{itemize}
\tightlist
\item
  KPI\#36: Guarantee 99.99999\% service availability and reliability
  based on effective monitoring and LCM.
\item
  KPI\#42: DT-validated AI algorithms that do not break the
  time-sensitive operation when deployed in 99\% of the outputs.
\item
  KPI\#37: Guarantee 99.99999\% service continuity through AI-driven
  proactive management of cloud-native services.
\end{itemize}}
\end{requirement}

\begin{requirement}{Req-T23-AE-05}{MUST}{Containment and automatic downgrade}
\reqfield{Description / rationale}{%
The platform shall not attempt to verify a publisher validity claim. An
evaluation function shall track outcome against declared objective per
agent, and sustained divergence beyond a declared auditable threshold
shall trigger automatic narrowing of scope and then suspension of
decision emission, enacted by the DMO as a registry authority change and
applied at the intent egress gate. The agent shall not be able to
decline it.}
\reqfield{Dependency}{%
Req-T23-AE-04}
\reqfield{Traceability (Forward)}{%
\begin{itemize}
\tightlist
\item
  Enacted by the DMO as a registry change.
\item
  Applied by the platform at the intent egress gate.
\end{itemize}}
\reqfield{UNITY-6G components}{%
\begin{itemize}
\tightlist
\item
  Evaluation function
\item
  DMO
\item
  Tuple registry
\end{itemize}}
\reqfield{Linked KPIs}{%
\begin{itemize}
\tightlist
\item
  KPI\#21: Reduce the performance penalty due to conflicts by 20\% by
  means of DLT, hierarchical algorithms and policy-based conflict
  resolution schemes.
\item
  KPI\#26: Increase the ratio of service management and LCM tasks
  resolved by local AE/DE components.
\item
  KPI\#42: DT-validated AI algorithms that do not break the
  time-sensitive operation when deployed in 99\% of the outputs.
\end{itemize}}
\end{requirement}

\begin{requirement}{Req-T23-AE-06}{MUST}{Separation of lifecycle and per-decision authority}
\reqfield{Description / rationale}{%
The evaluation function shall hold lifecycle authority only and shall
not veto individual decisions. Per-decision authority shall remain with
the conflict fabric. The two shall act on different objects and
different time scales.}
\reqfield{Dependency}{%
Req-T23-AE-05}
\reqfield{Traceability (Forward)}{%
\begin{itemize}
\tightlist
\item
  Structural: evaluation acts between cycles, the fabric within one.
\end{itemize}}
\reqfield{UNITY-6G components}{%
\begin{itemize}
\tightlist
\item
  Evaluation function
\item
  Conflict fabric
\end{itemize}}
\reqfield{Linked KPIs}{%
\begin{itemize}
\tightlist
\item
  KPI\#21: Reduce the performance penalty due to conflicts by 20\% by
  means of DLT, hierarchical algorithms and policy-based conflict
  resolution schemes.
\item
  KPI\#26: Increase the ratio of service management and LCM tasks
  resolved by local AE/DE components.
\end{itemize}}
\end{requirement}

\begin{requirement}{Req-T23-AE-07}{MUST}{Enumerated insight types enforced at the boundary}
\reqfield{Description / rationale}{%
A partial AE shall enumerate its insight types. The publication boundary
shall validate every insight against the registered enumeration and
schema; an undeclared type shall be dropped and shall raise a
conformance fault against the agent, and repeated faults shall feed the
downgrade path.}
\reqfield{Dependency}{%
Req-T23-FW-03}
\reqfield{Traceability (Forward)}{%
\begin{itemize}
\tightlist
\item
  Validated at the publication boundary.
\item
  Faults recorded against the agent.
\end{itemize}}
\reqfield{UNITY-6G components}{%
\begin{itemize}
\tightlist
\item
  AE runtime
\item
  Tuple registry
\end{itemize}}
\reqfield{Linked KPIs}{%
\begin{itemize}
\tightlist
\item
  KPI\#1: Deliver four distributed management entities (MS, AE, DE and
  ACT) deployable in each heterogeneous domain.
\item
  KPI\#26: Increase the ratio of service management and LCM tasks
  resolved by local AE/DE components.
\end{itemize}}
\end{requirement}

\begin{requirement}{Req-T23-AE-08}{MUST}{Mandatory insight provenance}
\reqfield{Description / rationale}{%
Every insight envelope shall carry the producing agent identifier and
version together with the model identifier and version. A consumer of
delegated analytics shall declare accepted provider ranges, and a
provider version outside that range shall raise a binding fault at the
consumer and trigger its declared model-unavailable behavior rather than
silently altering decision behavior.}
\reqfield{Dependency}{%
Req-T23-AE-01}
\reqfield{Traceability (Forward)}{%
\begin{itemize}
\tightlist
\item
  Populated by the AE runtime.
\item
  Range checked at the consumer binding.
\end{itemize}}
\reqfield{UNITY-6G components}{%
\begin{itemize}
\tightlist
\item
  AE runtime
\item
  Admission (IDMO/DMO)
\end{itemize}}
\reqfield{Linked KPIs}{%
\begin{itemize}
\tightlist
\item
  KPI\#42: DT-validated AI algorithms that do not break the
  time-sensitive operation when deployed in 99\% of the outputs.
\item
  KPI\#33: Support more tenants via distributed AI \& DLT.
\end{itemize}}
\end{requirement}

\begin{requirement}{Req-T23-AE-09}{MUST}{Model governance}
\reqfield{Description / rationale}{%
Each model shall be declared with identifier, version, required feature
set and version, validity conditions and serving mode, and shall be
bound by pinned identity and digest rather than by a mutable reference.
Model lifecycle operations shall never occur on the decision path.}
\reqfield{Dependency}{%
Req-T23-AE-04}
\reqfield{Traceability (Forward)}{%
\begin{itemize}
\tightlist
\item
  Checked by the validity rules at admission.
\end{itemize}}
\reqfield{UNITY-6G components}{%
\begin{itemize}
\tightlist
\item
  Tuple descriptor
\item
  Model registry
\end{itemize}}
\reqfield{Linked KPIs}{%
\begin{itemize}
\tightlist
\item
  KPI\#41: Reduction of 70\% of deployment time of AI-based
  resource/energy management algorithms.
\item
  KPI\#42: DT-validated AI algorithms that do not break the
  time-sensitive operation when deployed in 99\% of the outputs.
\end{itemize}}
\end{requirement}

\begin{requirement}{Req-T23-AE-10}{MUST}{Measured inference latency binds the effective class}
\reqfield{Description / rationale}{%
Inference latency shall be a mandatory observability metric emitted by
the AE runtime. Sustained breach of the effective latency class shall be
a conformance fault, shall feed the downgrade path, and shall cause the
effective class to be re-stamped downward in the registry so that
composition and priority logic use the measured figure.}
\reqfield{Dependency}{%
Req-T23-AE-05}
\reqfield{Traceability (Forward)}{%
\begin{itemize}
\tightlist
\item
  Emitted by the AE runtime.
\item
  Re-stamped in the registry by the DMO.
\end{itemize}}
\reqfield{UNITY-6G components}{%
\begin{itemize}
\tightlist
\item
  AE runtime
\item
  Tuple registry
\item
  DMO
\end{itemize}}
\reqfield{Linked KPIs}{%
\begin{itemize}
\tightlist
\item
  KPI\#4: Real-time reconfiguration of IEEE 802.11 network domain in
  range smaller than 10 ms from taking the decision.
\item
  KPI\#5: Real-time reconfiguration of cellular network domain in the
  range smaller than 10 ms from taking the decision.
\item
  KPI\#34: Reduce Open RAN real-time control loops time scales below 10
  ms and 10x in NTN for real-time 6G.
\end{itemize}}
\end{requirement}

\begin{requirement}{Req-T23-AE-11}{MUST}{Declared failure behavior and no undeclared fallback}
\reqfield{Description / rationale}{%
The AE failure policy shall declare responses to input staleness,
minimum input coverage, operation outside model validity, loss of a
serving endpoint and drift, emitting at minimum a drift signal. Silent
extrapolation shall be prohibited. A fallback heuristic shall be used
only where declared in advance; an undeclared fallback shall not be
used, and a declared fallback shall be flagged in the insight envelope.}
\reqfield{Dependency}{%
Req-T23-AE-02}
\reqfield{Traceability (Forward)}{%
\begin{itemize}
\tightlist
\item
  Checked by the validity rules.
\item
  Enforced by the AE runtime; fallback use flagged in the envelope.
\end{itemize}}
\reqfield{UNITY-6G components}{%
\begin{itemize}
\tightlist
\item
  Tuple descriptor
\item
  AE runtime
\end{itemize}}
\reqfield{Linked KPIs}{%
\begin{itemize}
\tightlist
\item
  KPI\#31: Reduce service recovery time below 180 s.
\item
  KPI\#36: Guarantee 99.99999\% service availability and reliability
  based on effective monitoring and LCM.
\item
  KPI\#37: Guarantee 99.99999\% service continuity through AI-driven
  proactive management of cloud-native services.
\end{itemize}}
\end{requirement}

\begin{requirement}{Req-T23-AE-12}{MUST}{Cross-domain analytics delegation}
\reqfield{Description / rationale}{%
An AE delegation may cross domain boundaries as a data product or an
exposed contract. It shall pin the provider interface profile and
accepted provider and model version ranges; admission shall verify that
the provider subject classes cover the consumer requirement and that
provider freshness meets consumer assumptions; and provider input
quality and validity state shall be forwarded within the insight. A
decision exchanged between domain agents is a proposal and shall confer
no authority over the receiving agent; the binding decision shall be
taken at the cross-domain management level.}
\reqfield{Dependency}{%
Req-T23-AE-08}
\reqfield{Traceability (Forward)}{%
\begin{itemize}
\tightlist
\item
  Scope and freshness checked at admission.
\item
  Quality and validity forwarded by the AE runtime.
\end{itemize}}
\reqfield{UNITY-6G components}{%
\begin{itemize}
\tightlist
\item
  Admission (IDMO/DMO)
\item
  AE runtime
\item
  Provider agent
\end{itemize}}
\reqfield{Linked KPIs}{%
\begin{itemize}
\tightlist
\item
  KPI\#26: Increase the ratio of service management and LCM tasks
  resolved by local AE/DE components.
\item
  KPI\#33: Support more tenants via distributed AI \& DLT.
\end{itemize}}
\end{requirement}

\hypertarget{contribution-to-the-project-kpis-1}{%
\subsection{CONTRIBUTION TO THE PROJECT KPIs}\label{contribution-to-the-project-kpis-1}}

Table 5-4 states what the Analytics Engine contributes to each KPI it
bears on, and Table 5-5 maps the requirements of this chapter onto those
KPIs.

\begin{wptable}{5-4}

\begin{longtable}[]{@{}
  >{\raggedright\arraybackslash}p{(\columnwidth - 4\tabcolsep) * \real{0.1499}}
  >{\raggedright\arraybackslash}p{(\columnwidth - 4\tabcolsep) * \real{0.6375}}
  >{\raggedright\arraybackslash}p{(\columnwidth - 4\tabcolsep) * \real{0.2126}}@{}}
\caption{Contribution of the Analytics Engine to the project
KPIs}\tabularnewline
\toprule\noalign{}
\begin{minipage}[b]{\linewidth}\raggedright
\textbf{KPI}
\end{minipage} & \begin{minipage}[b]{\linewidth}\raggedright
\textbf{Contribution of the Analytics Engine}
\end{minipage} & \begin{minipage}[b]{\linewidth}\raggedright
\textbf{Requirements}
\end{minipage} \\
\midrule\noalign{}
\endfirsthead
\toprule\noalign{}
\begin{minipage}[b]{\linewidth}\raggedright
\textbf{KPI}
\end{minipage} & \begin{minipage}[b]{\linewidth}\raggedright
\textbf{Contribution of the Analytics Engine}
\end{minipage} & \begin{minipage}[b]{\linewidth}\raggedright
\textbf{Requirements}
\end{minipage} \\
\midrule\noalign{}
\endhead
\bottomrule\noalign{}
\endlastfoot
\textbf{KPI\#1} & Supplies the AE of the four-entity tuple,
substitutable across domains because it holds no authority. &
Req-T23-AE-01, Req-T23-AE-07 \\
\textbf{KPI\#4} & Measured inference latency binds the effective class,
so a sub-10 ms loop is not claimed on a model that cannot serve within
it. & Req-T23-AE-10 \\
\textbf{KPI\#5} & The same measurement governs cellular loops. &
Req-T23-AE-10 \\
\textbf{KPI\#21} & Lifecycle authority and per-decision authority are
separated, so evaluation and conflict resolution never contend. &
Req-T23-AE-05, Req-T23-AE-06 \\
\textbf{KPI\#26} & Local analytics with local decision keeps resolution
inside the domain; cross-domain delegation remains available because
insight carries no authority. & Req-T23-AE-12 \\
\textbf{KPI\#31} & Declared responses to staleness, coverage and drift
mean a degraded model yields an unavailable-insight signal rather than a
late wrong answer. & Req-T23-AE-11 \\
\textbf{KPI\#33} & Cross-domain insight delegation supports more tenants
without granting authority across the boundary. & Req-T23-AE-12 \\
\textbf{KPI\#34} & Effective class re-stamping keeps real-time claims
tied to measured inference behavior. & Req-T23-AE-10 \\
\textbf{KPI\#36} & Degraded evidence propagates to the DE by
construction rather than by a model noticing its own degradation. &
Req-T23-AE-02, Req-T23-AE-03 \\
\textbf{KPI\#37} & Proactive management rests on forecasts whose
validity envelope is evaluated by the runtime before the forecast is
relied upon. & Req-T23-AE-04, Req-T23-AE-11 \\
\textbf{KPI\#41} & Pinned model identity and declared serving mode let
an algorithm be deployed without bespoke integration. & Req-T23-AE-09 \\
\textbf{KPI\#42} & Insight carries the provenance of agent and model,
input quality computed from the evidence consumed and a runtime-stamped
validity state, so out-of-envelope operation is stamped rather than
inferred; this is what allows a validated algorithm to be trusted not to
break time-sensitive operation. & Req-T23-AE-02, Req-T23-AE-04,
Req-T23-AE-05, Req-T23-AE-08 \\
\end{longtable}

\end{wptable}

\begin{wptable}{5-5}\footnotesize

\begin{longtable}[]{@{}
  >{\raggedright\arraybackslash}p{(\columnwidth - 24\tabcolsep) * \real{0.1814}}
  >{\raggedright\arraybackslash}p{(\columnwidth - 24\tabcolsep) * \real{0.0682}}
  >{\raggedright\arraybackslash}p{(\columnwidth - 24\tabcolsep) * \real{0.0682}}
  >{\raggedright\arraybackslash}p{(\columnwidth - 24\tabcolsep) * \real{0.0682}}
  >{\raggedright\arraybackslash}p{(\columnwidth - 24\tabcolsep) * \real{0.0682}}
  >{\raggedright\arraybackslash}p{(\columnwidth - 24\tabcolsep) * \real{0.0682}}
  >{\raggedright\arraybackslash}p{(\columnwidth - 24\tabcolsep) * \real{0.0682}}
  >{\raggedright\arraybackslash}p{(\columnwidth - 24\tabcolsep) * \real{0.0682}}
  >{\raggedright\arraybackslash}p{(\columnwidth - 24\tabcolsep) * \real{0.0682}}
  >{\raggedright\arraybackslash}p{(\columnwidth - 24\tabcolsep) * \real{0.0682}}
  >{\raggedright\arraybackslash}p{(\columnwidth - 24\tabcolsep) * \real{0.0682}}
  >{\raggedright\arraybackslash}p{(\columnwidth - 24\tabcolsep) * \real{0.0682}}
  >{\raggedright\arraybackslash}p{(\columnwidth - 24\tabcolsep) * \real{0.0679}}@{}}
\caption{Mapping of Analytics Engine requirements onto the project
KPIs}\tabularnewline
\toprule\noalign{}
\begin{minipage}[b]{\linewidth}\raggedright
\textbf{KPI\textbackslash{}}

\textbf{Requirement}
\end{minipage} & \begin{minipage}[b]{\linewidth}\raggedright
\textbf{\#1}
\end{minipage} & \begin{minipage}[b]{\linewidth}\raggedright
\textbf{\#4}
\end{minipage} & \begin{minipage}[b]{\linewidth}\raggedright
\textbf{\#5}
\end{minipage} & \begin{minipage}[b]{\linewidth}\raggedright
\textbf{\#21}
\end{minipage} & \begin{minipage}[b]{\linewidth}\raggedright
\textbf{\#26}
\end{minipage} & \begin{minipage}[b]{\linewidth}\raggedright
\textbf{\#31}
\end{minipage} & \begin{minipage}[b]{\linewidth}\raggedright
\textbf{\#33}
\end{minipage} & \begin{minipage}[b]{\linewidth}\raggedright
\textbf{\#34}
\end{minipage} & \begin{minipage}[b]{\linewidth}\raggedright
\textbf{\#36}
\end{minipage} & \begin{minipage}[b]{\linewidth}\raggedright
\textbf{\#37}
\end{minipage} & \begin{minipage}[b]{\linewidth}\raggedright
\textbf{\#41}
\end{minipage} & \begin{minipage}[b]{\linewidth}\raggedright
\textbf{\#42}
\end{minipage} \\
\midrule\noalign{}
\endfirsthead
\toprule\noalign{}
\begin{minipage}[b]{\linewidth}\raggedright
\textbf{KPI\textbackslash{}}

\textbf{Requirement}
\end{minipage} & \begin{minipage}[b]{\linewidth}\raggedright
\textbf{\#1}
\end{minipage} & \begin{minipage}[b]{\linewidth}\raggedright
\textbf{\#4}
\end{minipage} & \begin{minipage}[b]{\linewidth}\raggedright
\textbf{\#5}
\end{minipage} & \begin{minipage}[b]{\linewidth}\raggedright
\textbf{\#21}
\end{minipage} & \begin{minipage}[b]{\linewidth}\raggedright
\textbf{\#26}
\end{minipage} & \begin{minipage}[b]{\linewidth}\raggedright
\textbf{\#31}
\end{minipage} & \begin{minipage}[b]{\linewidth}\raggedright
\textbf{\#33}
\end{minipage} & \begin{minipage}[b]{\linewidth}\raggedright
\textbf{\#34}
\end{minipage} & \begin{minipage}[b]{\linewidth}\raggedright
\textbf{\#36}
\end{minipage} & \begin{minipage}[b]{\linewidth}\raggedright
\textbf{\#37}
\end{minipage} & \begin{minipage}[b]{\linewidth}\raggedright
\textbf{\#41}
\end{minipage} & \begin{minipage}[b]{\linewidth}\raggedright
\textbf{\#42}
\end{minipage} \\
\midrule\noalign{}
\endhead
\bottomrule\noalign{}
\endlastfoot
\textbf{Req-T23-AE-01} & \kpimark{} & & & & & & & & & & & \kpimark{} \\
\textbf{Req-T23-AE-02} & & & & & & & & & \kpimark{} & & & \kpimark{} \\
\textbf{Req-T23-AE-03} & & & & & & & & & \kpimark{} & & & \kpimark{} \\
\textbf{Req-T23-AE-04} & & & & & & & & & \kpimark{} & \kpimark{} & &
\kpimark{} \\
\textbf{Req-T23-AE-05} & & & & \kpimark{} & \kpimark{} & & & & & & &
\kpimark{} \\
\textbf{Req-T23-AE-06} & & & & \kpimark{} & \kpimark{} & & & & & & & \\
\textbf{Req-T23-AE-07} & \kpimark{} & & & & \kpimark{} & & & & & & & \\
\textbf{Req-T23-AE-08} & & & & & & & \kpimark{} & & & & & \kpimark{} \\
\textbf{Req-T23-AE-09} & & & & & & & & & & & \kpimark{} & \kpimark{} \\
\textbf{Req-T23-AE-10} & & \kpimark{} & \kpimark{} & & & & & \kpimark{}
& & & & \\
\textbf{Req-T23-AE-11} & & & & & & \kpimark{} & & & \kpimark{} &
\kpimark{} & & \\
\textbf{Req-T23-AE-12} & & & & & \kpimark{} & & \kpimark{} & & & & & \\
\end{longtable}

\end{wptable}
